\documentclass[11pt]{article}
\usepackage[margin=1in]{geometry}
\usepackage{amsmath,amssymb,booktabs,graphicx,xcolor,hyperref,natbib}
\usepackage[capitalise]{cleveref}

\newcommand{\todo}[1]{\textcolor{red}{[TODO: #1]}}
\newcommand{\pending}[1]{\textcolor{blue}{[PENDING: #1]}}
\newcommand{\ece}{\mathrm{ECE}}

\title{Region-Aware Gaussian Process Approximation Cannot Improve Calibration:\\
A Negative Result, Its Mechanism, and a Measurement Hazard}
\author{Srikar M.K.\\ Northeastern University\\ \texttt{manda.k@northeastern.edu}}
\date{\today}

\begin{document}
\maketitle

\begin{abstract}
A natural idea for scaling Gaussian processes is to spend approximation budget
unevenly: route ``important'' regions of the input space to high-fidelity
approximations and the rest to cheap ones. We show this cannot improve predictive
calibration, and identify why. The predictive variance of a sparse GP is dominated by
its shared global noise term rather than by the approximation gap, so models of
differing fidelity emit nearly identical predictive standard deviations at the same
input. On data whose regional structure is true by construction, the per-region
spread in calibration error across Random Fourier Features, Nystr\"om at two ranks,
and an exact GP is at most $0.009$ against a calibration error of $0.5$. Consequently
an oracle that selects the best model per region, using the \emph{true} generative
partition, improves expected calibration error by $+0.0003$ --- that is, not at all.
The same holds on five standard benchmarks under an oracle and a
proportion-preserving random-routing control.

We also report a measurement hazard encountered while establishing this. In a widely
used GP library, an opt-in setting documented as a fast variance approximation is
inert below a training-set size threshold that never appears in user code, and above
it perturbs predictive standard deviations by a median of ${\sim}100\%$, inflating
measured $\ece$ by $3$--$9\times$. Calibration results for exact GPs above this
threshold may not describe the model being reported.

Finally, we separate the idea's premise from its mechanism. Region structure does
matter, but the quantity that should vary by region is the \emph{noise model}, not
the approximation. \pending{final claim conditional on the heteroscedastic-GP
comparison in \cref{sec:noise}}
\end{abstract}

\section{Introduction}
\label{sec:intro}

Exact Gaussian process regression costs $O(n^3)$, and a large literature reduces this
with sparse approximations \citep{rasmussen2006, quinonero2005, titsias2009}. Most
such methods apply a single approximation uniformly across the input space. It is
natural to ask whether the budget should instead be allocated unevenly: regions that
are sparsely sampled, or where the model is uncertain, might warrant a more faithful
approximation than regions that are dense and easy.

We call this \emph{region-aware approximation}. The premise is appealing and the
mechanism is intuitive. This paper shows the mechanism cannot deliver on the premise,
for a reason that does not depend on the particular partitioning criterion, the
particular approximations, or the quality of the implementation.

\paragraph{Contributions.}
\begin{enumerate}
\item \textbf{An orthogonality result} (\cref{sec:orth}): approximation fidelity and
predictive uncertainty are nearly independent, because a sparse GP's predictive
variance is dominated by the shared noise term. We quantify this directly.
\item \textbf{A negative result with its ceiling} (\cref{sec:empirical}): even an
oracle given the true generative partition and the best model per region does not
beat using one model everywhere, on standard benchmarks and on data purpose-built to
favour routing.
\item \textbf{A measurement hazard} (\cref{sec:hazard}): approximate predictive
variance silently corrupts calibration metrics above a library-default threshold,
with no signal in user code.
\item \textbf{A redirection} (\cref{sec:noise}): the partition is useful, but for the
noise model rather than the approximation. \pending{strength of this claim}
\end{enumerate}

\section{Background}
\label{sec:background}

\paragraph{Sparse GP approximations.} Given data $\{(x_i,y_i)\}_{i=1}^n$ with
$y = f(x) + \varepsilon$, $\varepsilon \sim \mathcal{N}(0,\sigma_n^2)$, a GP prior
$f \sim \mathcal{GP}(0,k)$ yields a predictive distribution at $x_*$ with variance
\begin{equation}
\label{eq:var}
\mathbb{V}[y_*] \;=\; \underbrace{k(x_*,x_*) - k_*^\top (K+\sigma_n^2 I)^{-1} k_*}_{\text{latent, approximation-dependent}} \;+\; \underbrace{\sigma_n^2}_{\text{shared}} .
\end{equation}
Random Fourier Features \citep{rahimi2007} and Nystr\"om / subset-of-regressors
\citep{williams2001, smola2000} replace the first term with a rank-$m$ surrogate.
Crucially, they leave the second term untouched. \Cref{sec:orth} shows this is what
dooms region-aware approximation.

\paragraph{Calibration metrics.} For regression we use the expected calibration error
over central predictive intervals,
\begin{equation}
\ece = \frac{1}{B}\sum_{b=1}^{B} \Big| \tfrac{1}{n}\textstyle\sum_i \mathbf{1}\big[|y_i-\mu_i| \le z_{(1+c_b)/2}\,\sigma_i\big] - c_b \Big|,
\end{equation}
with $c_b$ spanning $[0.1,0.9]$, alongside Gaussian negative log-likelihood. Both
depend on $\sigma_i$, which is what makes \cref{sec:hazard} consequential.

\section{A measurement hazard in approximate predictive variance}
\label{sec:hazard}

Before reporting calibration comparisons we must establish that they measure the
model rather than the inference path.

GPyTorch \citep{gardner2018} provides \texttt{fast\_pred\_var()}, an implementation of
LOVE \citep{pleiss2018} that approximates the predictive covariance, and separately
falls back from Cholesky to iterative solves when the training set exceeds
\texttt{max\_cholesky\_size} (default $800$). We fit one model per cell and vary
\emph{only} the inference path.

\begin{table}[h]\centering
\caption{Same fitted model; only the inference path varies. Error is the median over
test points of $|\sigma_{\text{mode}} - \sigma_{\text{exact}}|/\sigma_{\text{exact}}$.
Onset is exact: over $20$ cells with $n\le800$ the maximum error is $0.0018\%$; over
$16$ cells with $n>800$ the minimum is $0.089\%$.}
\label{tab:hazard}
\begin{tabular}{lrrrr}
\toprule
dataset & $n_{\text{train}}$ & $\ece$ default & $\ece$ exact & median $\sigma$ error \\
\midrule
concrete   &  618 & 0.0434 & 0.0434 & \phantom{00}0.00\% \\
protein    & 4000 & 0.1084 & 0.0367 & \phantom{0}33.44\% \\
robot\_arm & 4000 & 0.2599 & 0.0824 & 199.22\% \\
sarcos     & 4000 & 0.3068 & 0.0302 & 213.62\% \\
synthetic  & 4000 & 0.4209 & 0.0654 & 559.12\% \\
\bottomrule
\end{tabular}
\end{table}

A four-way decomposition isolates the cause. Raising \texttt{max\_cholesky\_size}
alone leaves the error unchanged ($101.98\%$ median across cells with $n>800$);
disabling \texttt{fast\_pred\_var()} alone reduces it to $5.00\%$. The threshold
matters only because LOVE is a \emph{no-op} below it: the approximation is silently
inert on small data and silently severe on large data, with no change in user code.

\paragraph{Implication.} Reported $\ece$ or NLL for ``exact'' GPs trained above
$n=800$ in this library may describe an approximation. We recommend that calibration
studies disable fast predictive variance and raise the Cholesky threshold above $n$,
and report having done so.

\section{Fidelity and uncertainty are orthogonal}
\label{sec:orth}

\Cref{eq:var} suggests that if $\sigma_n^2$ dominates, then changing the rank of the
approximation --- which affects only the first term --- barely moves the predictive
standard deviation. We verify this directly on \texttt{hetero\_extreme}, a dataset
whose noise varies $50\times$ across three regions by construction.

\begin{table}[h]\centering
\caption{Per-region $\ece$ of four models of very different fidelity, evaluated at the
same inputs. Every model emits a mean predictive standard deviation of $0.563$ in all
three regions, while the true noise standard deviation varies from $0.020$ to $1.000$.}
\label{tab:orth}
\begin{tabular}{lrrrrrr}
\toprule
region & RFF & Nystr\"om-500 & Nystr\"om-1500 & exact GP & spread & true $\sigma_n$ \\
\midrule
0 & 0.4970 & 0.4979 & 0.4979 & 0.4982 & 0.0012 & 0.020 \\
1 & 0.3792 & 0.3791 & 0.3791 & 0.3751 & 0.0041 & 0.141 \\
2 & 0.1658 & 0.1648 & 0.1648 & 0.1740 & 0.0092 & 1.000 \\
\bottomrule
\end{tabular}
\end{table}

The spread across models \emph{within} a region is at most $0.009$, while the
calibration error itself is of order $0.5$. There is nothing to route between. This is
the paper's central mechanism: a router can only choose among the options it is given,
and along the axis that matters --- predictive uncertainty --- those options are
indistinguishable.

\section{Empirical study}
\label{sec:empirical}

\subsection{Protocol}
All methods share one split per (dataset, seed); hyperparameters are fitted by
marginal likelihood on training data; model capacity is selected on a validation
split; the test split is read once. A single $\ece$ estimator is used for every
method, and exact predictive variance is used throughout (\cref{sec:hazard}).

\subsection{Datasets are screened, not assumed}
\label{sec:screen}
Because the hypothesis concerns input-dependent structure, we measure it rather than
assume it. \Cref{tab:screen} reports $\mathrm{het}$, the ratio of the $90$th to $10$th
percentile of $k$NN-smoothed out-of-sample squared residuals.

\begin{table}[h]\centering
\caption{Measured input-dependent residual variance. Real data spans $2.4$--$7.1$.
Note that \texttt{synthetic\_heteroscedastic}, a dataset constructed specifically as a
controlled heteroscedastic test case, scores \emph{below} sarcos.}
\label{tab:screen}
\begin{tabular}{llr}
\toprule
group & dataset & het ratio \\
\midrule
constructed & hetero\_extreme & 1557.34 \\
constructed & varying\_smoothness & 22.33 \\
real & energy & 7.05 \\
standard & sarcos & 5.43 \\
real & airfoil & 4.44 \\
real & calhousing & 4.10 \\
standard & concrete & 3.88 \\
standard & synthetic\_heteroscedastic & 3.71 \\
real & yacht & 3.55 \\
standard & protein & 2.19 \\
\bottomrule
\end{tabular}
\end{table}

\subsection{Routing does not help, even given an oracle}
\label{sec:routing}

\begin{table}[h]\centering
\caption{$\ece$ on data whose regional structure is true by construction. \emph{true
oracle} uses the true generative partition and selects the best model per region on
validation data --- an upper bound on any region-aware router with these models.}
\label{tab:routing}
\begin{tabular}{lrrrrr}
\toprule
dataset & uniform best & true oracle & learned oracle & importance & random \\
\midrule
hetero\_extreme     & \textbf{0.2299} & 0.2315 & 0.2302 & 0.2322 & 0.2316 \\
varying\_smoothness & \textbf{0.0656} & 0.0658 & 0.0659 & 0.0661 & 0.0659 \\
varying\_density    & \textbf{0.0069} & 0.0075 & 0.0070 & 0.0069 & 0.0065 \\
combined            & 0.0525 & 0.0515 & \textbf{0.0504} & 0.0503 & 0.0537 \\
\bottomrule
\end{tabular}
\end{table}

The mean difference between the true oracle and a single model everywhere is
$+0.0003$ $\ece$. On five standard benchmarks the corresponding figure is $+0.0001$,
with the oracle losing on four of five. A proportion-preserving random-routing control
is statistically indistinguishable from importance-based routing.

\section{The partition is for the noise model}
\label{sec:noise}

\Cref{eq:var} also indicates the constructive move: the term that varies across
regions is $\sigma_n^2$, so that is what should be made regional. Holding the
partition, the posterior mean and the latent variance fixed, and replacing the scalar
noise with a per-region noise estimated from out-of-fold training residuals:

\pending{results table from \texttt{src/hetero\_fair.py}, 10 seeds, versus a
Kersting-style EM heteroscedastic GP with equal tuning budget. Prior expectation is
that a correctly implemented heteroscedastic GP matches or beats regional noise,
leaving the latter as a cheaper and more robust alternative rather than a better one.}

We additionally find that region identification must target the right quantity.
Scoring regions by predictive uncertainty and local sparsity --- the natural choice,
and the one used by the method that motivated this study --- recovers the true
partition only $41$--$49\%$ of the time against a chance rate of $33\%$. Scoring by
$k$NN-smoothed residual variance raises recovery to $96.9\%$ and $81.9\%$ on the two
datasets where noise regions exist, and correctly stays at chance on the dataset where
density varies but noise does not.

\section{Threats to validity}
\label{sec:threats}
\begin{itemize}
\item Our constructed datasets reach het ratios of ${\sim}1557$ where real data reaches
at most ${\sim}7$ (\cref{tab:screen}), and vary noise in steps rather than smoothly,
which structurally favours a piecewise-constant noise model. Claims in
\cref{sec:noise} rest on the real-data results, not the constructed ones.
\item The heteroscedastic baseline is our own implementation. \todo{add a comparison
against an independent implementation before submission}
\item Residual variance that varies with $x$ may reflect model bias rather than noise;
we therefore use the term \emph{input-dependent residual variance}.
\item \todo{$n$ capped at 2500--8000 throughout; scaling behaviour untested}
\end{itemize}

\section{Related work}
\todo{sparse GPs; heteroscedastic GPs \citep{goldberg1997, kersting2007}; calibration
in regression \citep{kuleshov2018}; locally adaptive and mixture-of-experts GPs, which
are the closest prior art and must be distinguished carefully}

\section{Conclusion}
Region-aware \emph{approximation} cannot improve calibration, because approximation
fidelity and predictive uncertainty are orthogonal. The premise behind it --- that
regions differ and a model should adapt --- is sound, but the adaptation must happen
in the noise model. Separately, practitioners measuring GP calibration should verify
that their inference path is exact.

\paragraph{Reproducibility.} All code, data and results:
\url{https://github.com/Srikarmk/Aurora-GP}. Every number above is regenerated by
\texttt{./reproduce.sh}.

\bibliographystyle{plainnat}
\bibliography{refs}
\end{document}
